\documentclass[12pt]{article}

\usepackage[a4paper,margin=1in]{geometry}
\usepackage{amsmath,amssymb,amsthm,bm,mathtools}
\usepackage{booktabs,threeparttable}
\usepackage{natbib}
\usepackage{hyperref}
\usepackage{enumitem}

\newtheorem{assumption}{Assumption}
\newtheorem{lemma}{Lemma}
\newtheorem{proposition}{Proposition}
\newtheorem{theorem}{Theorem}
\newtheorem{corollary}{Corollary}

\newcommand{\R}{\mathbb{R}}
\newcommand{\E}{\mathbb{E}}
\newcommand{\Var}{\operatorname{Var}}
\newcommand{\plim}{\operatorname*{plim}}
\newcommand{\op}{o_p}
\newcommand{\Op}{O_p}
\newcommand{\W}{W_N}
\newcommand{\SN}[1]{S_N(#1)}
\newcommand{\thetahat}{\widehat{\theta}}
\newcommand{\tauhat}{\widehat{\tau}}

\title{QML Estimation of a Spatial Panel Model with an Unknown Structural Break}
\author{Author Name}
\date{\today}

\begin{document}
\maketitle

\begin{abstract}
This paper studies quasi-maximum likelihood estimation for a spatial panel model with an unknown structural break. The abstract should summarize the model, estimator, asymptotic results, and Monte Carlo findings. This template contains placeholders that must be verified before submission.
\end{abstract}

\noindent\textbf{Keywords:} Spatial panel model; structural break; quasi-maximum likelihood; asymptotic theory; Monte Carlo simulation.\\
\textbf{JEL Codes:} C13; C21; C23.

\section{Introduction}\label{sec:introduction}

Spatial dependence and structural instability are two common features in panel data. Spatial econometric models capture interactions across cross-sectional units, while structural break models allow parameters to change across regimes. This paper develops a quasi-maximum likelihood framework for a spatial panel model with an unknown break date.

The key methodological challenge is that the likelihood contains both a spatial log-determinant term and a non-smooth break-date search. The proposed estimator profiles the likelihood over admissible break dates and estimates regime-specific spatial and slope parameters.

The paper makes three contributions. First, it specifies a spatial panel model with regime-specific parameters and an unknown common break date. Second, it develops a QML estimator and studies its asymptotic properties. Third, it evaluates finite-sample performance using Monte Carlo simulations.

\textbf{Placeholder warning.} The final contribution statement must be rewritten after completing the literature matrix.



\section{Model and Quasi-Likelihood}\label{sec:model-likelihood}

Let $i=1,\ldots,N$ index cross-sectional units and $t=1,\ldots,T$ index time. Let $y_t\in\R^N$ denote the dependent variable vector, $X_t\in\R^{N\times k}$ denote regressors, and $\W\in\R^{N\times N}$ denote a spatial weights matrix.

Consider the model
\begin{equation}
 y_t = \lambda_j \W y_t + X_t\beta_j + \W X_t\gamma_j + \alpha + u_t,
 \qquad t\in\mathcal{T}_j(\tau_0),\quad j=1,2.
 \label{eq:model-baseline}
\end{equation}
Equivalently,
\begin{equation}
 \SN{\lambda_j}y_t = X_t\beta_j + \W X_t\gamma_j + \alpha + u_t,
 \qquad \SN{\lambda_j}=I_N-\lambda_j\W.
 \label{eq:transformed-model}
\end{equation}

The true break date is $\tau_0$, with break fraction $\rho_0=\tau_0/T$. For a candidate break date $\tau$, define $\mathcal{T}_1(\tau)=\{1,\ldots,\tau\}$ and $\mathcal{T}_2(\tau)=\{\tau+1,\ldots,T\}$.

For $t\in\mathcal{T}_j(\tau)$, define residuals
\begin{equation}
 e_{jt}(\theta_j)=\SN{\lambda_j}y_t-X_t\beta_j-\W X_t\gamma_j-\alpha_j.
 \label{eq:residual}
\end{equation}

The Gaussian quasi-log-likelihood is
\begin{equation}
\ell_{NT}(\theta_1,\theta_2,\tau)
=\sum_{j=1}^2\left[
T_j(\tau)\log|\SN{\lambda_j}|
-\frac{NT_j(\tau)}{2}\log(2\pi\sigma_j^2)
-\frac{1}{2\sigma_j^2}\sum_{t\in\mathcal{T}_j(\tau)}e_{jt}(\theta_j)'e_{jt}(\theta_j)
\right].
\label{eq:loglik}
\end{equation}

The QML estimator is
\begin{equation}
(\widehat\theta_1,\widehat\theta_2,\widehat\tau)
=\arg\max_{\theta_1,\theta_2\in\Theta,\,\tau\in\mathcal{T}_\epsilon}
\ell_{NT}(\theta_1,\theta_2,\tau).
\label{eq:qml-estimator}
\end{equation}



\section{Assumptions}\label{sec:assumptions}

\begin{assumption}[Spatial weights]\label{ass:weights}
The spatial weights matrix $\W$ has zero diagonal, uniformly bounded row and column sums, and $\SN{\lambda}$ is nonsingular for all $\lambda\in\Lambda$.
\end{assumption}

\begin{assumption}[Parameter space]\label{ass:parameter-space}
The parameter space $\Theta$ is compact, and the true parameter lies in its interior. The variance parameters are bounded away from zero and infinity.
\end{assumption}

\begin{assumption}[Regressors and disturbances]\label{ass:regressors-errors}
The regressors and disturbances satisfy moment, dependence, and rank conditions sufficient for the uniform law of large numbers and central limit theorem used below.
\end{assumption}

\begin{assumption}[Break identification]\label{ass:break}
The true break fraction $\rho_0$ lies in $(\epsilon,1-\epsilon)$, and the pre-break and post-break regimes are distinguishable in the population objective.
\end{assumption}

\textbf{Placeholder warning.} Assumption~\ref{ass:regressors-errors} must be expanded into primitive conditions before the paper can be treated as theoretically complete.



\section{Asymptotic Properties}\label{sec:asymptotic-properties}

This section states candidate results. Each result must be verified against the final assumptions and proof details.

\begin{theorem}[Consistency]\label{thm:consistency}
Under Assumptions~\ref{ass:weights}--\ref{ass:break} and the uniform convergence conditions stated in Appendix~\ref{app:likelihood-expansion}, the QML estimator is consistent:
\[
\widehat\tau/T \to_p \rho_0,
\qquad
\widehat\theta_j \to_p \theta_{j0},\quad j=1,2.
\]
\end{theorem}

\begin{theorem}[Break-date rate: candidate]\label{thm:break-rate}
Under additional identification and stochastic equicontinuity conditions, the break-date estimator converges at a rate determined by the objective drift around $\tau_0$ and the likelihood fluctuation. The exact rate must be derived for the selected asymptotic regime.
\end{theorem}

\begin{theorem}[Asymptotic normality: candidate]\label{thm:asymptotic-normality}
Suppose the score satisfies a central limit theorem with effective information rate $A_{NT}$ and the Hessian converges to a nonsingular matrix. Then the regular parameter estimator satisfies
\[
A_{NT}^{1/2}(\widehat\theta-\theta_0)\Rightarrow \mathcal{N}(0,J^{-1}\Omega J^{-1}).
\]
\end{theorem}



\section{Monte Carlo Simulation}\label{sec:monte-carlo}

The Monte Carlo experiment evaluates the finite-sample performance of the QML estimator. For each design, the simulation reports bias, RMSE, empirical standard deviation, average estimated standard error, 95\% coverage probability, break-date accuracy, and convergence failure rate.

\begin{table}[htbp]
\centering
\caption{Monte Carlo results: placeholder}
\label{tab:mc-baseline}
\begin{threeparttable}
\begin{tabular}{rrrrrrrr}
\toprule
$N$ & $T$ & Parameter & True & Mean & Bias & RMSE & CP95 \\
\midrule
50 & 40 & $\lambda_1$ & 0.30 & -- & -- & -- & -- \\
50 & 40 & $\lambda_2$ & 0.50 & -- & -- & -- & -- \\
\bottomrule
\end{tabular}
\begin{tablenotes}
\small
\item Notes: This table is a placeholder. Final values must be generated from MATLAB outputs.
\end{tablenotes}
\end{threeparttable}
\end{table}



\section{Conclusion}\label{sec:conclusion}

This paper develops a QML framework for a spatial panel model with an unknown structural break. The final conclusion should summarize the verified theoretical results and simulation findings, and discuss extensions such as multiple breaks, dynamic spatial panels, heterogeneous break dates, or alternative spatial weight matrices.




\bibliographystyle{apalike}
\bibliography{refs}

\appendix
\section{Matrix Lemmas}\label{app:matrix-lemmas}

This appendix collects matrix results for $\W$ and $\SN{\lambda}$.

\begin{lemma}[Spatial multiplier bound: placeholder]\label{lem:spatial-multiplier-bound}
Under Assumption~\ref{ass:weights}, $\SN{\lambda}^{-1}$ is uniformly bounded for $\lambda\in\Lambda$.
\end{lemma}

\begin{proof}
Proof to be completed using the final spectral or norm condition on $\W$.
\end{proof}



\section{Likelihood Expansion}\label{app:likelihood-expansion}

This appendix derives the likelihood expansion and uniform convergence results.

\begin{lemma}[Uniform convergence: placeholder]\label{lem:uniform-convergence}
Under suitable moment and dependence conditions,
\[
\sup_{\theta,\tau}|Q_{NT}(\theta,\tau)-Q(\theta,\tau)|=o_p(1).
\]
\end{lemma}

\begin{proof}
Proof gap: specify primitive dependence and moment conditions and verify the uniform law of large numbers.
\end{proof}



\section{Proof of Consistency}\label{app:consistency-proof}

\begin{proof}[Proof of Theorem~\ref{thm:consistency}]
The proof follows from identification of the population objective and uniform convergence of the sample objective. Details must be completed after Lemma~\ref{lem:uniform-convergence} is proved.
\end{proof}



\section{Break-Date Rate}\label{app:break-rate-proof}

\begin{proof}[Proof of Theorem~\ref{thm:break-rate}]
Proof gap: derive the deterministic objective loss from assigning observations to the wrong regime and compare it with the stochastic fluctuation of the profile likelihood.
\end{proof}



\section{Asymptotic Normality}\label{app:asymptotic-normality-proof}

\begin{proof}[Proof of Theorem~\ref{thm:asymptotic-normality}]
Use a Taylor expansion of the score around the true parameter. Proof gap: verify the score CLT and Hessian convergence under the final assumptions.
\end{proof}




\end{document}


